% SPDX-License-Identifier: GPL-3.0-or-later
% Copyright (C) 2026 Alessandro Greco
\chapter{Архитектура сети}\label{cap:06-architettura}
\aggiornatoal{23 сентября 2026 г.}{tor 0.4.9.13; публичный консенсус от 23 сентября 2026 г., 15:00 UTC}

\begin{obiettivi}
К концу этой главы вы сможете:
\begin{itemize}[nosep]
  \item перечислять роли узлов сети Tor и их обязанности;
  \item описывать уровни протокола от TLS-канала до потока приложения;
  \item на высоком уровне прослеживать всё, что происходит при открытии веб-страницы через Tor;
  \item понимать, как разные реализации и версии сосуществуют в одной сети.
\end{itemize}
\end{obiettivi}

Прежде чем переходить к деталям следующих глав, полезно взглянуть на сеть сверху. Эта глава --- карта: позже мы вернёмся к каждому элементу и разберём его подробнее, а ссылки укажут соответствующие разделы.

% ---------------------------------------------------------------------------
\section{Роли}\label{sec:arch-ruoli}

Сеть Tor состоит из программ, выполняющих разные роли, зачастую несколько одновременно. \Cref{fig:arch-ruoli} сводит их вместе.

\begin{figure}[htbp]
\centering
% SPDX-License-Identifier: GPL-3.0-or-later
% Copyright (C) 2026 Alessandro Greco
%% Русское альтернативное описание EPUB.
\altfig{Схема ролей сети Tor. Вверху клиент достигает сайта через сторожевой, промежуточный и выходной ретрансляторы. Три значка условно представляют девять органов каталога, которые голосуют за сетевой консенсус и подписывают его. Для обхода цензуры мост может служить первым узлом. Внизу клиент и онион-сервис каждый строят свою цепь к узлу встречи без выходного ретранслятора; пунктирные линии не показывают промежуточные узлы.}
% Роли в сети Tor.
\begin{tikzpicture}[font=\sffamily\small,>=Stealth,
  aut/.style={relay,fill=torviola!15,draw=torviola,font=\sffamily\scriptsize},
  ex/.style={relay,fill=torverde!30}]
  \node[client] (c) at (0,0) {Клиент};
  \node[relay] (g) at (2.8,0) {сторожевой};
  \node[relay] (m) at (5.3,0) {средний};
  \node[ex] (e) at (7.8,0) {выходной};
  \node[server] (s) at (10.6,0) {Сайт};
  \draw[cellaf] (c) -- (g); \draw[cellaf] (g) -- (m); \draw[cellaf] (m) -- (e); \draw[->,thick,torgrigio] (e) -- (s);
  \foreach \x/\n in {3.8/1,5.3/2,6.8/3} \node[aut] (a\n) at (\x,2.4) {авт.};
  \node[font=\sffamily\scriptsize,above=1pt of a2] {органы каталога (9): подписанный консенсус};
  \draw[dotted,thick,torviola] (a2) -- node[font=\sffamily\scriptsize,right]{консенсус} (m);
  \node[relay,fill=torgrigio!10,draw=torgrigio,font=\sffamily\scriptsize] (b) at (1.4,-2.2) {мост};
  \node[font=\sffamily\scriptsize,align=center] at (1.4,-3.1) {первый узел вместо\\сторожевого при цензуре};
  \draw[dashed,torgrigio,->] (c) -- (b);
  \node[client,fill=torverde!15,draw=torverde!60!black] (o) at (8.8,-2.3) {Онион-сервис};
  \node[relay,fill=torverde!25,font=\sffamily\scriptsize,align=center] (rp) at (5.3,-2.3) {узел\\встречи};
  \draw[dashed,torverde!60!black,->] (o) -- node[font=\sffamily\scriptsize,below]{цепь} (rp);
  \draw[dashed,torverde!60!black,->] (c) to[bend right=18] node[font=\sffamily\scriptsize,above,pos=.68]{цепь клиента} (rp);
  \node[font=\sffamily\scriptsize,align=center] at (5.3,-3.35) {клиент и сервис встречаются\\без выхода через ретранслятор};
\end{tikzpicture}

\caption{Роли в сети Tor. Клиент строит цепь через guard-реле, промежуточное реле и exit-реле. Доверенные серверы каталогов (directory authorities) публикуют консенсус со списком реле. Для пользователей, сталкивающихся с цензурой, мосты заменяют guard-реле. Доступ к onion-сервисам выполняется через точку рандеву без использования exit-реле.}
\label{fig:arch-ruoli}
\end{figure}

\begin{description}
  \item[Клиент] Пользовательская программа, например входящая в Tor Browser. Она выбирает маршруты, строит цепи и предоставляет приложениям SOCKS-порт. Она не принимает входящие соединения из сети и не аутентифицирует себя (\cref{cap:07-canali}).
  \item[Реле] Узел, пересылающий чужой трафик. В зависимости от флагов, назначенных доверенными серверами каталогов, и своей конфигурации он может выполнять роль:
  \begin{itemize}
    \item \textbf{guard-реле}: первого перехода в цепях клиента, выбранного один раз и используемого длительное время (\cref{sec:guard});
    \item \textbf{промежуточного реле}: промежуточного перехода, не знающего ни источник, ни пункт назначения;
    \item \textbf{exit-реле}: последнего перехода, который открывает интернет-соединения согласно своей политике выхода (exit policy) (\cref{cap:11-stream});
    \item \textbf{HSDir}: хранилища дескрипторов onion-сервисов (\cref{sec:os-hsdir});
    \item \textbf{кэша каталогов}: узла, повторно распространяющего консенсус и дескрипторы (\cref{cap:12-directory}).
  \end{itemize}
  \item[Доверенный сервер каталогов] Один из девяти доверенных серверов, которые каждый час голосуют о состоянии сети и подписывают консенсус; некоторые также измеряют пропускную способность реле.
  \item[Мост] Непубличное реле, используемое первым переходом теми, кто не может обращаться к публичным реле, часто вместе с подключаемым транспортом (\cref{cap:16-censura}).
  \item[Onion-сервис] Сервис, доступный только через Tor без раскрытия своего местоположения (\cref{cap:14-onion-services}).
\end{description}

В публичном консенсусе от 23 сентября 2026 года насчитывается 9358 реле: 5306 с флагом \texttt{Guard} и 3177 с флагом \texttt{Exit}; 1615 имеют оба флага. На guard-реле приходится 70\,\% общей пропускной способности, на exit-реле --- 33\,\% (\cref{cap:13-percorso}).

% ---------------------------------------------------------------------------
\section{Уровни протокола}\label{sec:arch-livelli}

Протокол Tor можно представить как стек уровней, каждый из которых строится поверх нижележащего (\cref{fig:arch-livelli}).

\begin{figure}[htbp]
\centering
% SPDX-License-Identifier: GPL-3.0-or-later
% Copyright (C) 2026 Alessandro Greco
%% Русское альтернативное описание EPUB.
\altfig{Упрощённая схема шести уровней сверху вниз: приложение HTTP поверх потоков; команды потоков BEGIN, DATA, END, SENDME (глава 11); послойно зашифрованные сообщения RELAY (глава 10); управление цепью: CREATE2 и DESTROY — команды ячеек, а EXTEND2 — сообщение RELAY (глава 9); ячейки фиксированной длины 514 байт при версии протокола канала v4 и выше, наряду с ячейками переменной длины (глава 8); канал TLS поверх TCP с CERTS и NETINFO (глава 7). Это функциональные уровни, а не строгая классификация типов ячеек.}
% Уровни протокола Tor.
\begin{tikzpicture}[font=\sffamily\small,
  liv/.style={draw,rounded corners=2pt,minimum width=7cm,minimum height=0.8cm,align=center},
  rifcap/.style={font=\sffamily\scriptsize,text=torgrigio,anchor=west}]
  \node[liv,fill=torverde!20] (a) at (0,4) {Приложение (HTTP, \ldots) поверх потоков};
  \node[liv,fill=torverde!10] (s) at (0,3) {Поток: BEGIN, DATA, END, SENDME};
  \node[liv,fill=torviola!20] (r) at (0,2) {Послойно зашифрованные сообщения RELAY};
  \node[liv,fill=torviola!10] (c) at (0,1) {Цепи: CREATE2, EXTEND2, DESTROY};
  \node[liv,fill=torgrigio!15] (l) at (0,0) {Фикс. ячейки: 514 байт (v4+)};
  \node[liv,fill=torgrigio!5] (t) at (0,-1) {Канал: TLS поверх TCP, CERTS, NETINFO};
  \node[rifcap] at (3.7,3) {Глава 11};
  \node[rifcap] at (3.7,2) {Глава 10};
  \node[rifcap] at (3.7,1) {Глава 9};
  \node[rifcap] at (3.7,0) {Глава 8};
  \node[rifcap] at (3.7,-1) {Глава 7};
\end{tikzpicture}

\caption{Уровни протокола Tor и главы, в которых они описаны. Каждый уровень использует нижележащий: потоки передаются как сообщения RELAY, а сообщения RELAY помещаются в ячейки с послойным шифрованием, а ячейки идут внутри TLS-каналов.}
\label{fig:arch-livelli}
\end{figure}

\begin{description}
  \item[Канал] TLS-соединение между двумя соседними узлами, в рамках которого проверяется идентичность реле (\cref{cap:07-canali}).
  \item[Ячейки] Единицы передачи длиной 514 байт, переносимые каналом (\cref{cap:08-celle}).
  \item[Цепь] Многошаговый маршрут, строящийся по одному переходу с помощью протоколов согласования ключей, вырабатывающих разные ключи для каждого перехода (\cref{cap:09-circuiti}).
  \item[Послойное шифрование] Сообщения RELAY, обёрнутые одним слоем шифрования для каждого перехода (\cref{cap:10-cifratura-relay}).
  \item[Поток] TCP-соединения приложений, открываемые exit-реле и управляемые окнами и механизмом управления перегрузкой (\cref{cap:11-stream}).
\end{description}
Рядом с этим стеком существуют система каталогов, предоставляющая клиентам список реле, и протокол onion-сервисов, использующий цепи для встречи в точке рандеву.

% ---------------------------------------------------------------------------
\section{Веб-страница от начала до конца}\label{sec:arch-percorso}

Проследим на высоком уровне, что происходит, когда пользователь впервые запускает Tor Browser и открывает страницу.
\begin{enumerate}
  \item \textbf{Запуск.} tor связывается с доверенным сервером каталогов или одним из примерно двухсот \emph{резервных серверов каталогов} (fallback directories), перечисленных в исходном коде, и загружает консенсус, подписи которого проверяет. Затем загружаются микродескрипторы реле (\cref{cap:12-directory}).
  \item \textbf{Выбор guard-реле.} tor случайным образом выбирает небольшую выборку guard-реле с учётом пропускной способности и использует два из них (\cref{sec:guard}).
  \item \textbf{Канал.} tor открывает TLS-соединение с guard-реле и проверяет его идентичность с помощью ячеек \cella{CERTS} (\cref{cap:07-canali}).
  \item \textbf{Цепь.} tor выбирает exit-реле, подходящее для пункта назначения, и промежуточное реле, после чего строит трёхпереходную цепь с помощью \cella{CREATE2} и \texttt{EXTEND2} (\cref{cap:09-circuiti,cap:13-percorso}). На практике некоторые цепи строятся заранее, чтобы быть готовыми к моменту запроса.
  \item \textbf{Поток.} Через SOCKS-порт браузер запрашивает соединение с \texttt{www.example.com:443}. tor отправляет exit-реле сообщение \texttt{BEGIN}; exit-реле разрешает имя и открывает TCP-соединение (\cref{cap:11-stream}).
  \item \textbf{Данные.} Браузер общается с веб-сайтом по HTTPS через поток. Exit-реле видит имя сайта и TLS-трафик, но не знает, кто его запросил. Guard-реле видит, кто сделал запрос, но не знает, что именно было запрошено. Промежуточное реле не знает ни того, ни другого.
\end{enumerate}
Весь процесс запуска обычно занимает лишь несколько секунд. Многие механизмы последующих глав, включая заполнение, управление перегрузкой и vanguards, незаметно работают вокруг этой последовательности.

% ---------------------------------------------------------------------------
\section{Реализации и версии}\label{sec:arch-versioni}

\subsection{Две реализации}

Историческая реализация --- \textbf{tor} на C, которая выполняет все роли: клиента, реле, доверенного сервера каталогов и onion-сервиса. Наряду с ней Tor Project разрабатывает \textbf{Arti}, переписанную на Rust реализацию (\cref{cap:30-arti}), которая уже используется как клиент и для onion-сервисов. В публичном консенсусе от 23 сентября 2026 года Arti не используется ни одним реле: 6027 реле работают на tor 0.4.9.12, 2520 --- на 0.4.9.11, а большинство остальных --- на более старых выпусках серии 0.4.9. Сорок семь реле используют разрабатываемую версию 0.5.0.0-alpha-dev.

\subsection{Субпротоколы}

Как сеть из тысяч узлов, обновляющихся в разное время, вводит новые возможности? Tor не опирается для этого на номера версий программного обеспечения. Вместо этого используются \textbf{субпротоколы}: каждая возможность относится к субпротоколу с собственным номером версии, а каждое реле объявляет поддерживаемые версии в строке \texttt{pr} консенсуса~\spec{tor-spec}{Subprotocol-based versioning}. Например, реле в публичной сети объявляет:
\begin{lstlisting}[numbers=none]
pr Conflux=1 Cons=1-2 Desc=1-4 DirCache=2 FlowCtrl=1-2 HSDir=2 HSIntro=4-5
   HSRend=1-2 Link=3-5 LinkAuth=3 Microdesc=1-3 Padding=2 Relay=2-6
\end{lstlisting}
Клиент, желающий использовать ntor-v3 с реле, проверяет, объявляет ли оно \texttt{Relay=4}; CGO требует \texttt{Relay=6}, а управление перегрузкой --- \texttt{FlowCtrl=2}. Эти числа обозначают независимые \emph{возможности}, а не накопительные версии: формально \texttt{Relay=6} не подразумевает \texttt{Relay=4} (\proposta{346}, поэтому там предпочитается термин \enquote{возможности субпротокола}). На практике tor объявляет их непрерывным диапазоном, поскольку реализует функции по порядку. В консенсусе также перечислены \emph{обязательные} субпротоколы: версия tor, не поддерживающая их, понимает, что устарела, и отказывается работать, вместо того чтобы подвергать пользователей риску.

% ---------------------------------------------------------------------------
\section*{Итоги}
\begin{itemize}
  \item Сеть состоит из клиентов, реле (guard, middle, exit, HSDir и кэшей каталогов), доверенных серверов каталогов, мостов и onion-сервисов.
  \item Протокол образует стек: TLS-каналы, ячейки, цепи, послойное шифрование и потоки.
  \item Guard-реле знает пользователя, но не пункт назначения; exit-реле знает пункт назначения, но не пользователя.
  \item tor (C) и Arti (Rust) сосуществуют. Возможности согласуются как субпротоколы, объявляемые в строке \texttt{pr} консенсуса.
\end{itemize}

% ---------------------------------------------------------------------------
\section*{Упражнения}

\begin{esercizio}[1]\label{es:06-vista}
Для каждой роли --- guard, middle, exit, HSDir и directory authority --- укажите, что она знает о пользователе, посещающем веб-сайт через Tor.
\end{esercizio}

\begin{esercizio}[1]\label{es:06-pr}
Какие функции, рассматриваемые в главах 9--11, поддерживает реле из строки \texttt{pr}, приведённой в этой главе? Можно ли использовать его как exit-реле с CGO?
\end{esercizio}

\begin{esercizio}[2]\label{es:06-livelli}
Через Tor загружается веб-страница размером 100\,KB. Оцените по уровням, сколько байт проходит по линии между клиентом и guard-реле, учитывая ячейки, заголовки сообщений RELAY и SENDME.
\end{esercizio}

% ---------------------------------------------------------------------------
\section*{Дополнительная литература}
\begin{itemize}
  \item R.~Dingledine, N.~Mathewson, P.~Syverson, \titolo{Tor: The Second-Generation Onion Router}~\cite{dingledine2004tor}: исходная архитектура, значительная часть которой применяется до сих пор.
  \item Tor Metrics, \url{https://metrics.torproject.org/}: актуальные сведения о реле, пользователях и версиях.
\end{itemize}
